\paragraph{A cohomological obstruction to naive cuspidal devissage.}
There is a concrete loss of information when one tries to
iterate maximal cuspidally pro-\(\ell\) reconstruction through
ramified covers. Work over an algebraically closed field of
characteristic two, or over \(\mathbb C\), and take the
pro-odd groups throughout. Let \(C\) have genus \(g\geq2\),
let \(P\in C\), and write
\[
 F=\pi_1(C-P)^{\mathrm{odd}},\quad
 G=\pi_1(C)^{\mathrm{odd}},\quad K=\ker(F\to G),
 \quad N=\ker(K\to K^{(3)}).
\]
Here \(K^{(3)}\) denotes the maximal pro-three quotient,
so \(F/N\) is the maximal cuspidally pro-three quotient over
\(G\). Choose a cyclic etale cubic \(a:C'\to C\), and a
cyclic cubic \(D\to C'\) branched exactly at the three
points of \(a^{-1}(P)\), with ramification index three at
each. Such a cover exists because the degree-three branch
divisor has a cube root in \(\operatorname{Pic}(C')\):
choose a line bundle of degree one and correct its cube
by divisibility of \(\operatorname{Pic}^0(C')\) over the
algebraically closed field.
It has genus \(h=9g-5\).

Put \(D^\circ=D-\pi^{-1}(P)\). Its degree-nine etale cover
of \(C-P\) is visible in \(F/N\), since its Galois closure
has three-group Galois group. Let \(H\leq F/N\) be its
open subgroup, and let \(\overline H\) be the quotient of
\(H\) by the closed normal subgroup generated by its three
puncture inertia groups. Then, for every odd prime
\(\ell\ne3\),
\[
 \dim_{\mathbb F_\ell}H^1(\overline H,\mathbb F_\ell)
       =2g(C')=6g-4
       <18g-10
       =\dim_{\mathbb F_\ell}H^1(D,\mathbb F_\ell).
                                                               \tag{AG.39.11}
\]
Thus \(\overline H\) is a proper quotient of the full
group \(\pi_1(D)^{\mathrm{odd}}\), already detected on
finite abelian quotients. For \(g=2\), the two dimensions
are \(8\) and \(26\).

Here is a direct proof of the comparison, including its
surjectivity assertion. Let \(F'\leq F\) correspond to
\(C'-a^{-1}(P)\), and let
\(\chi:F'\twoheadrightarrow\mathbb Z/3\) define
\(D^\circ\). The group \(K\) is contained in \(F'\),
and \(\chi|_K\) is surjective, because a branch inertia
generator lies in \(K\) and maps nontrivially. If
\(I=\ker\chi\), this gives \(F'=IK\). Also \(N\subset I\),
and hence
\[
 1\longrightarrow (I\cap K)/N\longrightarrow
 H=I/N\longrightarrow \pi_1(C')^{\mathrm{odd}}
 \longrightarrow1.                                  \tag{AG.39.12}
\]
The left-hand term is a closed subgroup of the pro-three
group \(K/N\). Every homomorphism from \(H\) to
\(\mathbb F_\ell\) therefore factors uniquely through
the last group when \(\ell\ne3\). The puncture inertia
groups in \(H\) lie in the left-hand term, so killing them
does not change these homomorphisms. This proves the first
equality in (AG.39.11); the second is the usual surface
cohomology dimension and the computed genus of \(D\).
The natural map from \(\pi_1(D)^{\mathrm{odd}}\) to
\(\overline H\) is surjective by the construction.

This is an obstruction to the specific procedure of taking
an open subgroup of a cuspidally pro-three quotient and then
killing its puncture inertia. It is not a theorem that all
cuspidal quotients together determine insufficient information
by every possible method. In a solvable induction, the next
cuspidally pro-\(\ell\) theorem needs the \emph{full} proper
prime-to-two group of the intermediate ramified compactification.
Equation (AG.39.11) shows exactly why the preceding formal
operations do not provide that input. That ramified
compactification is also absent from the original proper
unramified tower.
